\documentclass[11pt,a4paper]{article}

% ------------------------------------------------------------------
%  Rapport de Lab Cyber -- Mentorat Cybersecurite / FORCE N
%  Lab : DHCP Starvation sur routeur Cisco (GNS3 + c7200)
% ------------------------------------------------------------------

\usepackage[utf8]{inputenc}
\usepackage[T1]{fontenc}
\usepackage[french]{babel}
\usepackage[a4paper,margin=2.3cm]{geometry}
\usepackage{xcolor}
\usepackage{graphicx}
\usepackage{listings}
\usepackage{fancyhdr}
\usepackage{titlesec}
\usepackage{enumitem}
\usepackage{array}
\usepackage{booktabs}
\usepackage{longtable}
\usepackage{tcolorbox}
\usepackage[hidelinks]{hyperref}

% ----- Couleurs de la charte -----
\definecolor{forceblue}{HTML}{1E70BF}
\definecolor{forcegreen}{HTML}{3AAA35}
\definecolor{boxbg}{HTML}{F4F4F4}
\definecolor{warnred}{HTML}{C0392B}

% ----- Titres en bleu -----
\titleformat{\section}{\normalfont\large\bfseries\color{forceblue}}{\thesection.}{0.6em}{}
\titleformat{\subsection}{\normalfont\normalsize\bfseries\color{forceblue}}{\thesubsection}{0.6em}{}
\titleformat{\subsubsection}{\normalfont\normalsize\bfseries}{}{0em}{}

% ----- En-tete / pied de page -----
\pagestyle{fancy}
\fancyhf{}
\setlength{\headheight}{30pt}
\lhead{\raisebox{-0.3\height}{\IfFileExists{images/logo_forcen.png}{\includegraphics[height=0.85cm]{images/logo_forcen.png}}{\footnotesize\color{forceblue}FORCE N}}}
\rhead{\small\color{forceblue} DHCP Starvation (Cisco / GNS3), Mentorat Cybersécurité}
\cfoot{\thepage}
\renewcommand{\headrulewidth}{0.4pt}

% ----- Style des encadres de commandes -----
\lstdefinestyle{cmd}{
  basicstyle=\ttfamily\small,
  frame=single,
  rulecolor=\color{black!55},
  backgroundcolor=\color{boxbg},
  breaklines=true,
  columns=fullflexible,
  keepspaces=true,
  xleftmargin=4pt,
  framexleftmargin=4pt,
  showstringspaces=false,
  literate={é}{{\'e}}1 {è}{{\`e}}1 {à}{{\`a}}1 {ê}{{\^e}}1 {î}{{\^i}}1 {ô}{{\^o}}1 {ç}{{\c c}}1 {É}{{\'E}}1
}
\lstset{style=cmd}

% ----- Macro d'insertion de capture (placeholder si le fichier manque) -----
\newcommand{\capture}[2]{%
  \begin{center}
  \IfFileExists{#1}{\includegraphics[width=0.88\linewidth]{#1}}{%
    \setlength{\fboxsep}{10pt}%
    \fbox{\parbox[c][3.2cm][c]{0.85\linewidth}{\centering\itshape%
      [ Capture à insérer ]\\ \texttt{\detokenize{#1}}}}}\\[3pt]
  {\small\itshape #2}
  \end{center}}

% ----- Encadre d'avertissement -----
\newtcolorbox{avertissement}{colback=warnred!6,colframe=warnred,boxrule=0.8pt,arc=2pt,
  fonttitle=\bfseries,title=Avertissement légal et éthique}

% ----- Bloc signature (auteur + empreinte GPG) -----
% =====================================================================
%  Bloc signature / marque — HAMADY CISSE
%  Réutilisable : copier ce fichier dans un projet et faire
%      % =====================================================================
%  Bloc signature / marque — HAMADY CISSE
%  Réutilisable : copier ce fichier dans un projet et faire
%      % =====================================================================
%  Bloc signature / marque — HAMADY CISSE
%  Réutilisable : copier ce fichier dans un projet et faire
%      \input{signature.tex}      (dans le préambule)
%  puis appeler  \signatureHamady  là où la signature doit apparaître.
%  Dépendances : xcolor (déjà chargé dans les rapports du mentorat).
% =====================================================================

\definecolor{sigblue}{HTML}{1E70BF}

% Empreinte GPG (fingerprint) — à remplacer après génération de la clé.
\newcommand{\gpgEmpreinte}{A6F4 0DFD EE3E B7AE 80B3\ 6E77 361C 9E02 B790 303E}

\newcommand{\signatureHamady}{%
  \par\vspace{1.2em}%
  \noindent{\color{sigblue}\rule{\linewidth}{0.6pt}}\par\vspace{0.4em}%
  \noindent{\bfseries\large HAMADY CISSE}\par\vspace{0.15em}%
  \noindent{\small Diplômé en Cryptographie et Sécurité de l'Information}\par\vspace{0.35em}%
  \noindent{\footnotesize\ttfamily Courriel : hamadycisse7540@gmail.com \quad GitHub : hamadycisse7540-hash}\par\vspace{0.1em}%
  \noindent{\footnotesize\ttfamily Empreinte GPG : \gpgEmpreinte}\par\vspace{0.3em}%
  \noindent{\scriptsize\itshape Document signé numériquement (GPG). Vérifiez son authenticité et son intégrité avec la clé publique ci-dessus.}\par\vspace{0.4em}%
  \noindent{\color{sigblue}\rule{\linewidth}{0.6pt}}\par
}
      (dans le préambule)
%  puis appeler  \signatureHamady  là où la signature doit apparaître.
%  Dépendances : xcolor (déjà chargé dans les rapports du mentorat).
% =====================================================================

\definecolor{sigblue}{HTML}{1E70BF}

% Empreinte GPG (fingerprint) — à remplacer après génération de la clé.
\newcommand{\gpgEmpreinte}{A6F4 0DFD EE3E B7AE 80B3\ 6E77 361C 9E02 B790 303E}

\newcommand{\signatureHamady}{%
  \par\vspace{1.2em}%
  \noindent{\color{sigblue}\rule{\linewidth}{0.6pt}}\par\vspace{0.4em}%
  \noindent{\bfseries\large HAMADY CISSE}\par\vspace{0.15em}%
  \noindent{\small Diplômé en Cryptographie et Sécurité de l'Information}\par\vspace{0.35em}%
  \noindent{\footnotesize\ttfamily Courriel : hamadycisse7540@gmail.com \quad GitHub : hamadycisse7540-hash}\par\vspace{0.1em}%
  \noindent{\footnotesize\ttfamily Empreinte GPG : \gpgEmpreinte}\par\vspace{0.3em}%
  \noindent{\scriptsize\itshape Document signé numériquement (GPG). Vérifiez son authenticité et son intégrité avec la clé publique ci-dessus.}\par\vspace{0.4em}%
  \noindent{\color{sigblue}\rule{\linewidth}{0.6pt}}\par
}
      (dans le préambule)
%  puis appeler  \signatureHamady  là où la signature doit apparaître.
%  Dépendances : xcolor (déjà chargé dans les rapports du mentorat).
% =====================================================================

\definecolor{sigblue}{HTML}{1E70BF}

% Empreinte GPG (fingerprint) — à remplacer après génération de la clé.
\newcommand{\gpgEmpreinte}{A6F4 0DFD EE3E B7AE 80B3\ 6E77 361C 9E02 B790 303E}

\newcommand{\signatureHamady}{%
  \par\vspace{1.2em}%
  \noindent{\color{sigblue}\rule{\linewidth}{0.6pt}}\par\vspace{0.4em}%
  \noindent{\bfseries\large HAMADY CISSE}\par\vspace{0.15em}%
  \noindent{\small Diplômé en Cryptographie et Sécurité de l'Information}\par\vspace{0.35em}%
  \noindent{\footnotesize\ttfamily Courriel : hamadycisse7540@gmail.com \quad GitHub : hamadycisse7540-hash}\par\vspace{0.1em}%
  \noindent{\footnotesize\ttfamily Empreinte GPG : \gpgEmpreinte}\par\vspace{0.3em}%
  \noindent{\scriptsize\itshape Document signé numériquement (GPG). Vérifiez son authenticité et son intégrité avec la clé publique ci-dessus.}\par\vspace{0.4em}%
  \noindent{\color{sigblue}\rule{\linewidth}{0.6pt}}\par
}


\begin{document}

% ================================================================
%                        PAGE DE TITRE
% ================================================================
\begin{titlepage}
\centering
\vspace*{1.5cm}

\IfFileExists{images/logo_forcen.png}{\includegraphics[width=7cm]{images/logo_forcen.png}}{\fbox{\parbox[c][3cm][c]{7cm}{\centering\itshape [ logo\_forcen.png ]}}}\par
\vspace{1.2cm}

{\Large\bfseries\color{forceblue} Mentorat CYBERSÉCURITÉ CYB6 / Groupe 3\par}
\vspace{0.4cm}
{\large Encadrant : M. Ibrahima Samba DIOP\par}
{\small Consultant en Cybersécurité\par}
\vspace{2.2cm}

\rule{\linewidth}{0.5pt}\\[0.4cm]
{\huge\bfseries Rapport de Lab\par}
\vspace{0.4cm}
{\LARGE DHCP Starvation sur routeur Cisco\\ (GNS3 + IOS c7200)\par}
\rule{\linewidth}{0.5pt}\\[1.5cm]

\begin{tabular}{rl}
\textbf{Étudiant(e) :} & HAMADY CISSE \\[0.5cm]
\textbf{Date :} & \today \\[0.5cm]
\textbf{Environnement :} & GNS3 + Cisco c7200, réseau \texttt{192.168.6.0/24} \\
\end{tabular}

\vfill
{\small Document réalisé dans un cadre strictement pédagogique et sur un laboratoire isolé.}
\end{titlepage}

\tableofcontents
\newpage

% ================================================================
\section{Introduction}
% ================================================================
Ce rapport reprend l'attaque \textbf{DHCP Starvation} mais la reproduit sur un
\textbf{véritable routeur Cisco}, dans l'environnement d'origine prévu par l'encadrant : une
topologie \textbf{GNS3} dans laquelle un routeur Cisco joue le rôle de serveur DHCP. Il
complète le rapport précédent, où le serveur DHCP était une machine VyOS sous VirtualBox.

L'intérêt de cette seconde version est double : d'une part, coller à la maquette du cours
(routeur Cisco « RouteurDHCP ») ; d'autre part, observer une \textbf{différence de comportement}
entre deux serveurs DHCP face à la même attaque. Le routeur est émulé à partir d'une image IOS
\texttt{c7200-adventerprisek9-mz.153-3.XB12} fournie par l'encadrant.

\begin{avertissement}
Toutes les manipulations décrites ont été réalisées exclusivement sur un laboratoire
\textbf{isolé} (topologie GNS3 locale, sans exposition vers un réseau de production ni vers
Internet), composé d'équipements émulés et de machines virtuelles appartenant à l'auteur.
N'effectuez jamais ce type d'attaque sur un site, un système ou un réseau sans l'autorisation
préalable et écrite du propriétaire.
\end{avertissement}

\subsection{Objectifs}
\begin{itemize}[leftmargin=1.4em]
  \item Monter sous GNS3 la topologie du cours avec un \textbf{vrai routeur Cisco} (image IOS
        c7200) comme serveur DHCP.
  \item Rattacher une machine \textbf{Kali} (VirtualBox) à la topologie et vérifier qu'elle
        reçoit une adresse du routeur Cisco.
  \item Réaliser l'attaque DHCP Starvation avec \texttt{yersinia} et prouver l'épuisement du
        pool via \texttt{show ip dhcp binding}.
  \item Comparer le comportement du serveur Cisco à celui d'un serveur VyOS/Kea face à la même
        attaque.
\end{itemize}

% ================================================================
\section{Environnement du laboratoire}
% ================================================================
La topologie est montée sous \textbf{GNS3} : un routeur Cisco \textbf{c7200} (nommé \texttt{R1})
sert de serveur DHCP, relié par un \textbf{switch Ethernet} à une machine \textbf{Kali Linux}
(VirtualBox) qui joue l'attaquant, sur le réseau \texttt{192.168.6.0/24}.

\begin{center}
\renewcommand{\arraystretch}{1.3}
\begin{tabular}{>{\bfseries}l l l l}
\toprule
Équipement & Rôle & Adresse IP & Interface / MAC \\
\midrule
R1 (Cisco c7200) & Serveur DHCP / passerelle & \texttt{192.168.6.5}         & \texttt{FastEthernet0/0} \\
Switch1          & Switch (Ethernet switch GNS3) & \texttt{---}             & \texttt{---} \\
vm-kali          & Attaquant (client DHCP)   & \texttt{192.168.6.11} (DHCP) & \texttt{08:00:27:aa:ac:ee} \\
\bottomrule
\end{tabular}
\end{center}

\noindent\emph{Note : le document de l'encadrant s'appuyait sur VMware Workstation et son
réseau NAT (VMnet8). Ici, la topologie est entièrement interne à GNS3 (routeur, switch et VM
reliés dans GNS3) : l'étape d'identification du réseau NAT VMware ne s'applique donc pas, mais
le plan d'adressage reste identique (\texttt{192.168.6.0/24}).}

\subsection{Topologie GNS3}
\capture{images/00_topologie_gns3.png}{Topologie sous GNS3 : \texttt{R1} (routeur Cisco c7200, serveur DHCP) $\rightarrow$ \texttt{Switch1} $\rightarrow$ \texttt{vm-kali-1} (attaquant). Tous les liens sont actifs.}

% ================================================================
\section{Tâche 1 : Mise en place de la topologie GNS3}
% ================================================================
Avant de configurer le DHCP, il faut un routeur Cisco fonctionnel dans GNS3 et une Kali reliée
à la topologie.

\subsection{Rappels théoriques}
GNS3 émule les routeurs Cisco à partir d'une \textbf{image IOS} via \textbf{Dynamips}. Un
routeur c7200 n'a pas besoin de virtualisation matérielle (il n'entre donc pas en conflit avec
VirtualBox). Une machine VirtualBox existante peut être intégrée comme nœud de la topologie.

\subsection{Manipulations}
\textbf{1. Importer l'image IOS du routeur :} dans \emph{Edit $\rightarrow$ Preferences
$\rightarrow$ Dynamips $\rightarrow$ IOS routers}, on crée un modèle à partir du fichier
\texttt{c7200-adventerprisek9-mz.153-3.XB12.image}, et on calcule une valeur \emph{Idle-PC}
(pour éviter 100\% de CPU).
\capture{images/01_c7200_import.png}{Modèle \texttt{c7200} créé dans GNS3 : image IOS chargée, Idle-PC calculé (\texttt{0x60630d08}), slot 0 = C7200-IO-FE (interface FastEthernet0/0).}

\textbf{2. Rattacher la machine Kali (VirtualBox) :} dans \emph{Preferences $\rightarrow$
VirtualBox VMs}, on ajoute la VM \texttt{vm-kali}. Point important : dans l'onglet
\emph{Network} du nœud, il faut \textbf{décocher} « Allow GNS3 to use any configured VirtualBox
adapter » pour que GNS3 câble réellement l'adaptateur de Kali au switch de la topologie (sinon
Kali reste sur son réseau VirtualBox d'origine et n'obtient pas d'adresse).
\capture{images/01_vmkali_template.png}{Ajout de la VM \texttt{vm-kali} comme nœud GNS3.}

\textbf{3. Câbler la topologie :} \texttt{R1 (FastEthernet0/0)} $\rightarrow$ \texttt{Switch1}
$\rightarrow$ \texttt{vm-kali}, puis démarrer les nœuds. Les liens passent au vert lorsque tout
est actif.

\subsection{Questions de la tâche 1}
\textbf{Q1. Pourquoi l'émulation du c7200 ne crée-t-elle pas de conflit avec VirtualBox,
contrairement à ce qu'on redoutait ?}

\textit{Réponse :} Le c7200 est émulé par \textbf{Dynamips}, qui reproduit le processeur du
routeur de façon logicielle et n'utilise pas la virtualisation matérielle (VT-x). Or le
conflit classique sous Windows oppose VirtualBox à l'hyperviseur (Hyper-V), tous deux
demandeurs de VT-x. Comme Dynamips ne réclame pas VT-x, GNS3 (pour le routeur) et VirtualBox
(pour Kali) cohabitent sans se disputer la virtualisation matérielle.

% ================================================================
\section{Tâche 2 : Configuration du serveur DHCP sur le routeur Cisco}
% ================================================================
On configure R1 exactement comme le « RouteurDHCP » du document de l'encadrant.

\subsection{Manipulations}
\textbf{1. Configurer l'interface et le pool DHCP :}
\begin{lstlisting}
enable
configure terminal
 interface fastEthernet 0/0
  ip address 192.168.6.5 255.255.255.0
  no shutdown
 exit
 ip dhcp pool L3SIMAC
  network 192.168.6.0 255.255.255.0
  default-router 192.168.6.2
 exit
 ip dhcp excluded-address 192.168.6.2
 ip dhcp excluded-address 192.168.6.255
 end
write memory
\end{lstlisting}
\capture{images/02_r1_config.png}{Configuration de R1 : interface FastEthernet0/0 en 192.168.6.5, pool L3SIMAC, adresses exclues. \texttt{show ip interface brief} confirme l'interface \texttt{up/up}.}

\textbf{2. Vérifier que Kali reçoit une adresse du routeur Cisco :} une fois Kali démarrée, son
interface \texttt{eth0} reçoit automatiquement une adresse du pool.
\begin{lstlisting}
ip a
ping -c 3 192.168.6.5
\end{lstlisting}
\capture{images/02_kali_dhcp.png}{Kali a reçu \texttt{192.168.6.11/24} du routeur Cisco (R1). Le \texttt{ping} vers R1 répond avec \texttt{ttl=255}, signature d'un équipement Cisco.}

\textbf{3. Tester la connectivité dans l'autre sens (routeur $\rightarrow$ Kali) :} depuis la
console de R1, on vérifie qu'on joint bien le client (comme dans le document de l'encadrant).
\begin{lstlisting}
ping 192.168.6.11
\end{lstlisting}
\capture{images/02_ping_r1_kali.png}{Depuis R1, \texttt{ping 192.168.6.11} : \texttt{Success rate is 100 percent (5/5)}. La connectivité est bidirectionnelle.}

\subsection{Questions de la tâche 2}
\textbf{Q1. Le \texttt{ttl=255} de la réponse ping vers 192.168.6.5 : qu'indique-t-il ?}

\textit{Réponse :} La valeur de TTL de départ dépend du système. Les équipements Cisco (IOS)
initialisent le TTL à 255, là où un hôte Linux part en général de 64 et Windows de 128. Un
\texttt{ping} qui revient avec \texttt{ttl=255} confirme donc que c'est bien le routeur Cisco
(et non un autre hôte) qui répond, ce qui valide que Kali dialogue avec le bon serveur.

% ================================================================
\section{Tâche 3 : Installation de yersinia et Wireshark (Partie B)}
% ================================================================
L'attaque s'appuie sur \texttt{yersinia} (attaques de niveau 2) et sur \textbf{Wireshark} pour
l'analyse du trafic. Sur la Kali du lab, yersinia avait déjà été installé selon la procédure de
l'encadrant. Pour mémoire, l'installation depuis les sources se fait ainsi :
\begin{lstlisting}
sudo git clone https://github.com/tomac/yersinia /opt/yersinia
sudo apt install autoconf libgtk-3-dev libnet1-dev libgtk2.0-dev libpcap-dev -y
cd /opt/yersinia
sudo ./autogen.sh
sudo ./configure --with-gtk
sudo make
sudo make install
\end{lstlisting}

\noindent Sur une Kali récente, la compilation des sources échoue avec le compilateur actuel.
Le paquet officiel s'installe alors directement :
\begin{lstlisting}
sudo apt install -y yersinia
\end{lstlisting}
Ce paquet étant compilé sans interface graphique (\texttt{-{}-disable-gtk}), on utilise
l'interface texte \texttt{sudo yersinia -I} au lieu de \texttt{yersinia -G}. Wireshark est
présent par défaut sur Kali (\texttt{sudo wireshark}).

% ================================================================
\section{Tâche 4 : Attaque DHCP Starvation sur le routeur Cisco}
% ================================================================
On sature maintenant le serveur DHCP du routeur Cisco et on observe la table de baux se remplir.

\subsection{Manipulations}
\textbf{1. État initial du serveur (avant attaque) :}
\begin{lstlisting}
show ip dhcp binding
\end{lstlisting}
\capture{images/03_binding_avant.png}{Avant l'attaque, la table ne contient qu'un bail : \texttt{192.168.6.11} (Kali, Client-ID \texttt{0108.0027.aaac.ee}).}

\textbf{2. Lancer le flood de DISCOVER avec yersinia (sur Kali) :} dans \texttt{yersinia -I},
mode DHCP (touche \texttt{g}), puis attaque \texttt{1} (\emph{DOS sending DISCOVER packet},
touche \texttt{x}).
\capture{images/03_yersinia_flood.png}{yersinia en mode DHCP : plus de 155\,000 paquets DISCOVER envoyés depuis \texttt{0.0.0.0} vers \texttt{255.255.255.255}, avec \texttt{MAC Spoofing [X]} actif (une MAC source différente par paquet).}

\textbf{3. Constater l'épuisement du pool (sur R1) :}
\begin{lstlisting}
show ip dhcp binding
\end{lstlisting}
\capture{images/03_binding_apres.png}{La table de baux du routeur Cisco est désormais pleine : les adresses du pool (192.168.6.1, .3, .4, \dots) sont attribuées à des MAC falsifiées. Le pool est épuisé.}

\textbf{4. Analyser le trafic avec Wireshark (sur Kali) :}
\begin{lstlisting}
sudo wireshark
\end{lstlisting}
\capture{images/03_wireshark.png}{Capture Wireshark (filtre \texttt{dhcp}) sur \texttt{eth0} : flot continu de \texttt{DHCP Discover} émis depuis \texttt{0.0.0.0} vers \texttt{255.255.255.255}, avec des adresses MAC source usurpées. Signature typique du DHCP Starvation.}

\subsection{Questions de la tâche 4}
\textbf{Q1. Pourquoi, sur le routeur Cisco, un simple flood de DISCOVER (yersinia) suffit-il à
remplir la table de baux, alors qu'il ne suffisait pas sur VyOS/Kea ?}

\textit{Réponse :} C'est la principale observation de ce lab. Le serveur DHCP de Cisco IOS
inscrit une entrée de type \emph{Automatic} dès qu'il attribue une adresse à une requête, ce
qui fait apparaître les adresses offertes dans \texttt{show ip dhcp binding} au fil du flood.
Le moteur \textbf{Kea} de VyOS, lui, ne valide un bail qu'au terme du cycle DORA complet
(DISCOVER $\rightarrow$ OFFER $\rightarrow$ REQUEST $\rightarrow$ ACK) : un flood de DISCOVER
seuls y laissait la table vide, et il fallait compléter le handshake (script Scapy) pour vider
le pool. Même attaque, deux implémentations : le résultat visible dépend du serveur.

\medskip\noindent\textbf{Q2. En quoi l'usurpation de MAC (MAC Spoofing) est-elle le cœur de
l'attaque ?}

\textit{Réponse :} Le serveur identifie chaque client par son adresse MAC et réserve une
adresse par MAC. En changeant la MAC source à chaque paquet, yersinia simule des centaines de
clients distincts et force le routeur à puiser dans son pool jusqu'à l'épuiser. Sans usurpation,
le serveur n'attribuerait qu'une seule adresse.

% ================================================================
\section{Réflexion globale et conclusion}
% ================================================================
\begin{enumerate}[leftmargin=1.6em]
  \item \textbf{Difficultés rencontrées et solutions :}
  \begin{itemize}[leftmargin=1.2em]
    \item \emph{Modèle Kali fantôme dans GNS3} : un modèle \texttt{kali-linux-2022.1} pointait
          vers une VM inexistante (erreur \texttt{VBOX\_E\_OBJECT\_NOT\_FOUND}) $\rightarrow$
          utiliser la vraie VM \texttt{vm-kali}.
    \item \emph{Kali sans adresse, lien rouge dans GNS3} : l'adaptateur n'était pas câblé au
          switch $\rightarrow$ \textbf{décocher} « Allow GNS3 to use any configured VirtualBox
          adapter » et démarrer la VM depuis GNS3 (VM éteinte au préalable).
  \end{itemize}

  \item \textbf{Comparaison Cisco vs VyOS/Kea :} la même attaque donne un résultat observable
        différent selon le serveur. Sur Cisco, le flood de DISCOVER remplit directement la table
        de baux ; sur Kea, il faut le cycle DORA complet. Cela rappelle qu'une attaque doit
        toujours être évaluée au regard de l'implémentation ciblée.

  \item \textbf{Mesures de défense :} comme pour la version VyOS, la parade se situe sur le
        \textbf{switch d'accès} : \textbf{DHCP snooping} (ports de confiance, limite de débit,
        filtrage des réponses pirates) et \textbf{port-security} (limite du nombre de MAC par
        port, qui coupe net l'usurpation de MAC). Ces contre-mesures font l'objet d'une
        démonstration dédiée.
\end{enumerate}

\medskip
\noindent\textbf{Conclusion.} Cette seconde version confirme l'attaque sur un véritable routeur
Cisco, dans la maquette du cours : quelques dizaines de secondes de \texttt{yersinia} suffisent
à épuiser le pool, et \texttt{show ip dhcp binding} se remplit de baux fantômes. La leçon reste
la même que pour la version VyOS : un protocole sans authentification du client se retourne
facilement en déni de service, et seul un durcissement de la couche d'accès (DHCP snooping,
port-security) protège réellement le réseau.

\signatureHamady

% ================================================================
\section{Annexe : commandes utilisées, rôles et options}
% ================================================================

\renewcommand{\arraystretch}{1.25}
\begin{longtable}{@{}p{6.3cm} p{8.5cm}@{}}
\toprule
\textbf{Commande} & \textbf{Rôle et options} \\
\midrule
\endfirsthead
\toprule
\textbf{Commande} & \textbf{Rôle et options} \\
\midrule
\endhead
\midrule
\multicolumn{2}{r}{\footnotesize\itshape suite page suivante} \\
\endfoot
\bottomrule
\endlastfoot

\multicolumn{2}{@{}l}{\textbf{\color{forceblue}Configuration du serveur DHCP (Cisco IOS)}} \\*
\midrule
\lstinline!enable! & Passe en mode privilégié (EXEC). \\
\lstinline!configure terminal! & Entre en configuration globale. \\
\lstinline!interface fastEthernet 0/0! & Sélectionne l'interface du serveur. \\
\lstinline!ip address 192.168.6.5 255.255.255.0! & Affecte l'IP et le masque /24. \\
\lstinline!no shutdown! & Active l'interface. \\
\lstinline!ip dhcp pool L3SIMAC! & Crée le pool DHCP \texttt{L3SIMAC}. \\
\lstinline!network 192.168.6.0 255.255.255.0! & Plage distribuable du pool. \\
\lstinline!default-router 192.168.6.2! & Passerelle annoncée aux clients. \\
\lstinline!ip dhcp excluded-address 192.168.6.2! & Retire la passerelle du pool. \\
\lstinline!ip dhcp excluded-address 192.168.6.255! & Retire l'adresse de broadcast. \\
\lstinline!write memory! & Enregistre la configuration en NVRAM. \\
\lstinline!show ip interface brief! & Affiche l'état et l'IP des interfaces. \\
\lstinline!show ip dhcp binding! & Liste les baux DHCP : sert à constater l'épuisement du pool. \\
\midrule

\multicolumn{2}{@{}l}{\textbf{\color{forceblue}Réseau et attaque (Kali)}} \\*
\midrule
\lstinline!ip a! & Affiche les interfaces et adresses de Kali. \\
\lstinline!ping -c 3 192.168.6.5! & Teste la connectivité vers R1 ; \texttt{-c 3} limite à 3 paquets. \\
\lstinline!sudo yersinia -I! & Interface texte de yersinia ; dans l'outil : \texttt{g} (DHCP) puis \texttt{x} + \texttt{1} (flood DISCOVER). \\

\end{longtable}

\end{document}
